\documentclass[11pt,a4paper]{article}
\usepackage[margin=27mm]{geometry}
\usepackage{fontspec,amsmath,amssymb,amsthm,unicode-math}
\setmainfont{STIXGeneral}\setmathfont{STIX Two Math}
\setlength{\emergencystretch}{2em}
\newcommand{\Sym}{\operatorname{Sym}}\newcommand{\Alt}{\operatorname{Alt}}
\newcommand{\gen}[1]{\langle#1\rangle}\newcommand{\eps}{\varepsilon}
\newtheorem{lemma}{Lemma}\newtheorem{theorem}[lemma]{Theorem}
\title{A transitive input to the rank and degree extension}
\author{Human Compression project}\date{Private bounded correction; distinct review pending}
\begin{document}\maketitle
\begin{abstract}
We construct a rank-five string C-group acting as $S_{11}$, with a perfect
split satisfying the hypotheses of the rank and degree extension in the
specified Cameron--Fernandes--Leemans author edition. Its prescribed extension
fails the intersection property. The input C-property is proved through
explicit subgroup formulas and the standard end criterion. Thus no finite
census is a premise of the counterexample. The example is below the rank
range of the source's main stability theorem.
\end{abstract}

\section{The marking and the source claim}
On $\{1,\ldots,11\}$ define
\[
\begin{aligned}
 x&=(8\ 9)(10\ 11),&y&=(1\ 8)(9\ 10),\\
 q&=(1\ 2)(3\ 4),&r&=(2\ 6)(3\ 5)(4\ 7),
\quad s&=(3\ 4)(6\ 7).
\end{aligned}
\]
The ordered marking is $(x,y,q,r,s)$. All five permutations are distinct
nonidentity involutions; distant generators commute. Composition below uses
right actions, so $z^{ab}=(z^a)^b$ and $a^b=b^{-1}ab$.
For a supported permutation $g$, write $\eps(g)\in\mathbb F_2$ for its parity.

Our exact source is Cameron--Fernandes--Leemans,
\emph{The number of string C-groups of high rank},
arXiv:2212.12723v2, 11 January 2023 (L019).
Its Theorem 6.4, pp.~32--33, concerns a C-marking with a perfect split,
a fracture graph, both split endpoints moved by the deleted subgroup,
and one trivial internal factor of the split generator.
It asserts that replacing that generator by the prescribed two generators
at a new point preserves the C-property. We verify the printed definitions
on pp.~8--9 and 16--17 directly after proving that this input is C.
We retain the standard end-parabolic criterion, source Proposition 2.1,
attributed there to McMullen--Schulte, Proposition 2E16. Every subgroup
formula used with that criterion is proved below.

\section{A small permutation calculation and a parity character}
Put $t=(3\ 4)$, $u=(3\ 5)(4\ 7)$, $D=\gen{t,u}$ on
$Y=\{3,4,5,7\}$, and $\tau=(9\ 10)$.
The product $tu$ is a $4$-cycle, so $D$ is dihedral of order eight.
There is a homomorphism
$\chi:D\longrightarrow\mathbb F_2$ with $\chi(t)=\chi(u)=1$:
its kernel is the cyclic rotation subgroup $\gen{tu}$.
This is a character of the dihedral marking, not necessarily the ordinary
parity of its action on $Y$.

\begin{lemma}\label{lem:Q}
The marking $(q,r,s)$ is a rank-three C-marking of
$\Sym(\{1,2,3,4,5,6,7\})$.
\end{lemma}
\begin{proof}
Direct multiplication gives the short word
\[
 w=qr(qsr)^4s=(4\ 5).
\]
For clarity, the following complete table lists the conjugates needed:
\[
\begin{array}{c|rrrrrr}
z&1&q&r&rs&rsr&rsrq\\\hline
w^z&(4\ 5)&(3\ 5)&(3\ 7)&(4\ 6)&(2\ 7)&(1\ 7).
\end{array}
\]
Each equality is checked on the seven displayed points. The six edges in
the bottom row form a spanning tree. Transpositions on the edges of a
connected graph generate the whole symmetric group: along a path one
obtains its endpoint transposition by successive conjugation. Hence
$\gen{q,r,s}=S_7$.

On $\{1,2,6\}$ the generators $q,r$ act as $(1\ 2),(2\ 6)$;
on $Y$ they act as $t,u$. Consequently
\begin{equation}\label{eq:T}
 T:=\gen{q,r}
 =\{g d:g\in\Sym(\{1,2,6\}),\ d\in D,\ \eps(g)=\chi(d)\}.
\end{equation}
Indeed the group is contained in this set, of order $24$, while $qr$
has order $\operatorname{lcm}(3,4)=12$. Two involutions whose product
has order $12$ generate a dihedral group of order $24$, proving equality.
The product $rs$ is the cycle $(2\ 7\ 3\ 5\ 4\ 6)$, so
$R:=\gen{r,s}$ is the dihedral group of order $12$ in its natural
six-point action. An element of $T\cap R$ fixes $1$ and preserves
$\{1,2,6\}$, hence preserves the adjacent pair $\{2,6\}$ of this
six-cycle. Its setwise stabilizer in the dihedral group is
$\gen{r}$, of order two: no nonidentity rotation preserves an edge,
and precisely one reflection interchanges its endpoints.
Thus $T\cap R=\gen{r}$. The two rank-two end markings are C, and
the standard end criterion proves the lemma.
\end{proof}

We will also use the elementary following observation. If a transitive
permutation group $L$ on a set $V$ fixes a new point $b$, then adjoining
$(a\ b)$ for $a\in V$ generates $\Sym(V\cup\{b\})$.
Conjugation by $L$ gives every $(v\ b)$, and these generate that
symmetric group. This observation is applied to restriction images below;
it does not assume that the corresponding supported transpositions have
already been separated in the original group.

\section{The first rank-four end}
Let
\[
 A=\gen{x,y,q,r},\quad N=\gen{x,y,q},\quad M=\gen{y,q,r},
 \quad P=\gen{y,q}.
\]
The two generators $x,y$ form the natural dihedral group of order ten on
$V=\{1,8,9,10,11\}$: their product is the $5$-cycle
$(1\ 8\ 10\ 11\ 9)$. They are even permutations.
The projection of $N$ onto $V\cup\{2\}$ is therefore $S_6$, by adjoining
the transposition $(1\ 2)$ as in the preceding observation. Since the
only action on $Y$ is the factor $t$ of $q$, one obtains
\begin{equation}\label{eq:N}
 N=\{g t^{\eps(g)}:g\in\Sym(\{1,2,8,9,10,11\})\}.
\end{equation}
The parity condition is respected by each generator; surjectivity of the
projection and uniqueness of its permitted second coordinate prove the formula.
Similarly,
\begin{equation}\label{eq:P}
 P=\{g t^e\tau^v:g\in\Sym(\{1,2,8\}),\ \eps(g)=e+v\}.
\end{equation}
The projection is $S_3$ and $(yq)^3=t\tau$, so it contains both points
of each permitted fiber. Its order is $12$.
The group $P$ fixes $11$, while the stabilizer of $11$ in the natural
$D_{10}=\gen{x,y}$ is $\gen{y}$. Thus
$\gen{x,y}\cap P=\gen{y}$, and the end criterion makes $N$ a C-marking.

Next put $Z=\{1,2,6,8\}$. Then
\begin{equation}\label{eq:M}
 M=\{g d\tau^v:g\in\Sym(Z),\ d\in D,\ \eps(g)=\chi(d)+v\}.
\end{equation}
To prove surjectivity onto $\Sym(Z)\times D$, omit the $\tau$ coordinate.
The images of $y,q,r$ on $Z$ are the three adjacent transpositions on
the path $8,1,2,6$, so their projection is $S_4$. Moreover $y$ projects
to a pure transposition in this product, with trivial $D$ coordinate.
Conjugating it by elements projecting to all of $S_4$ yields every
transposition in $S_4\times1$. Cancelling these first coordinates then
recovers both $t,u$ in the second factor. The image is the full product.
The displayed parity equation holds for all three original generators
and uniquely determines $v$, proving \eqref{eq:M}.

The end subgroups of this marking are $P$ and $T$ from \eqref{eq:T}.
A common element has its root permutation supported on $\{1,2\}$.
Membership in $T$ forces $v=0$, while membership in $P$ forces its
$D$ coordinate to be $t^e$ with $e=\eps(g)$. Hence
$P\cap T=\gen{q}$, proving that $M$ is C.

The full first end has formula
\begin{equation}\label{eq:A}
 A=\{g d:g\in\Sym(X),\ d\in D,\ \eps(g)=\chi(d)\},
 \qquad X=\{1,2,6,8,9,10,11\}.
\end{equation}
Its projection onto $X$ is $S_7$: the preceding $S_6$ projection becomes
$S_7$ on adjoining the image $(2\ 6)$ of $r$. The parity equation
holds for all generators. In addition, $(qr)^4$ is a pure $3$-cycle
on $\{1,2,6\}$, since its $D$ coordinate is $(tu)^4=1$.
Its conjugates under the surjective $S_7$ projection are every $3$-cycle
on $X$, with trivial $D$ coordinate. They generate $\Alt(X)$.
Since the projection onto $D$ is onto, multiplying any lift of $d$ by
these even permutations gives every root permutation of the required
parity. This proves \eqref{eq:A}, without an assumption about unknown
normal subgroups.

Finally \eqref{eq:N} and \eqref{eq:M} show $N\cap M=P$.
Indeed a common permutation fixes $6$ and $11$, has an arbitrary root
permutation $g$ on $\{1,2,8\}$, a possible factor $\tau^v$ on
$\{9,10\}$, and a $D$ coordinate restricted to $t^e$.
Both formulas then give the same equation $\eps(g)=e+v$,
which is exactly \eqref{eq:P}. Thus both end markings of $A$ are C
and meet in the required middle group, so $A$ is C.

\section{The second end and the whole input}
Put $B=\gen{y,q,r,s}$ and $Q=\gen{q,r,s}=S_7$.
The projection of $B$ onto $\{1,\ldots,8\}$ is $S_8$.
Since $(1\ 2)\in Q$ and
\[
 (y(1\ 2))^3=\tau,
\]
we obtain
\begin{equation}\label{eq:B}
 B=\Sym(\{1,\ldots,8\})\times\gen{\tau}.
\end{equation}
The end markings of $B$ are $M$ and $Q$, both already proved C.
Their intersection fixes $8,9,10,11$, so \eqref{eq:M} reduces to
$g\in\Sym(\{1,2,6\})$, $v=0$ and $\eps(g)=\chi(d)$.
This is precisely $T=\gen{q,r}$. Therefore $B$ is C.

For $A\cap B$, formula \eqref{eq:B} forces a permutation in the
$X$ coordinate of \eqref{eq:A} to preserve $Z=\{1,2,6,8\}$,
preserve $\{9,10\}$, and fix $11$. Write it as $g\tau^v$ with
$g\in\Sym(Z)$. Its parity equation in \eqref{eq:A} becomes
$\eps(g)+v=\chi(d)$, exactly the equation defining $M$.
Thus $A\cap B=M$. The end criterion proves that
$(x,y,q,r,s)$ is a string C-marking.

Its full group is $S_{11}$, not merely a transitive subgroup.
Formula \eqref{eq:B} supplies the supported $S_8$ and $\tau=(9\ 10)$.
Conjugating $(1\ 8)\in S_8$ by $x$ gives $(1\ 9)$, which yields
$S_9$. Adjoining $\tau$ yields $S_{10}$. Cancelling $(8\ 9)$ from
$x=(8\ 9)(10\ 11)$ now yields $(10\ 11)$ and hence $S_{11}$.

\section{The split, the fracture graph and the failed extension}
For deletion of labels $0,1,2,3,4$, respectively, the orbit partitions are
\[
\begin{array}{c|l}
0&\{1,2,3,4,5,6,7,8\},\{9,10\},\{11\}\\
1&\{1,2,3,4,5,6,7\},\{8,9\},\{10,11\}\\
2&\{1,8,9,10,11\},\{2,3,4,5,6,7\}\\
3&\{1,2,8,9,10,11\},\{3,4\},\{5\},\{6,7\}\\
4&\{1,2,6,8,9,10,11\},\{3,4,5,7\}.
\end{array}
\]
These partitions follow either from the formulas above or directly from
connectivity in the deleted permutation graphs. Each has more than the
one orbit of the full group. The respective edges
$\{8,9\},\{1,8\},\{1,2\},\{2,6\},\{6,7\}$
join distinct orbits in their rows, and form a path with isolated vertices
added. Thus they are a fracture graph in the printed definition.

Take the perfect split at label $i=2$, with
$a=1$, $b=2$, $A_0=\{1,8,9,10,11\}$ and
$B_0=\{2,3,4,5,6,7\}$. Deleting $q$ gives exactly these two orbits.
The unique crossing edge is $\{1,2\}$; the other edge $\{3,4\}$ lies
inside $B_0$. Both earlier marks are supported on $A_0$, both later
marks on $B_0$, and the split factorization is
$q=(1\ 2)\beta$ with $\alpha=1$ and $\beta=(3\ 4)$.
The two endpoints move under $y$ and $r$, respectively.
All hypotheses of source Theorem 6.4 are now verified, including transitivity.

At the new point $c=12$ the prescribed marking is
\[
 (x,y,u_0,v_0,r,s),\qquad
 u_0=(1\ 12),\quad v_0=(2\ 12)(3\ 4).
\]
It satisfies the string commutations. However
$(u_0v_0)^3=(3\ 4)=\beta$.
The group $\gen{v_0,r,s}$ is the copy of the $S_7$ in
Lemma~\ref{lem:Q} obtained by renaming point $1$ as $12$.
Consequently it contains $\beta$. Equivalently the same word $w$, with
$q$ replaced by $v_0$, gives $(4\ 5)$; conjugation by $v_0$ gives
$(3\ 5)$, and these give $(3\ 4)$.
But $\beta\notin\gen{v_0}$, since it is nonidentity and fixes $12$,
whereas $v_0$ moves $12$.
Thus the subgroups on labels $\{2,3\}$ and $\{3,4,5\}$ meet in
more than the subgroup on their common label $\{3\}$.
The new marking is not C.

\section{Limits of the conclusion}
This refutes the printed Theorem 6.4 even for transitive inputs.
It does not refute the main stability theorem: this example has rank
$r=5$ and degree $n=11$, so it does not meet the stability range
$r\geq(n+3)/2=7$. No conclusion about that remaining high-rank case,
a corrected sufficient condition, full L019 recovery or historical
priority is asserted here. The separate frozen disconnected example and
the valid forward endpoint extension used in L010 remain unchanged.
The computations supplied with this note are diagnostics of its explicit
permutation and subgroup proof, not a substitute for the written argument.
\end{document}
